\section{Two elementary analytic tools}
\label{sec:analytic-tools}

We record the Fourier and gamma-function facts used in the primitive
functional equation. Their usual normalizations can be compared with
\cite[\S\S1.8(iv), 5.2]{DLMF}.

\subsection{Poisson summation and the Gaussian}

A function $f:\mathbb R\to\mathbb C$ is a Schwartz function if it is
smooth and $\sup_x |x^j f^{(k)}(x)|<\infty$ for every pair of
nonnegative integers $j,k$. Define its Fourier transform by
\[
 \widehat f(y)=\int_{\mathbb R}f(x)e^{-2\pi ixy}\,dx.
\]
\begin{lemma}\label{tools:poisson}
For a Schwartz function,
\[
 \sum_{n\in\mathbb Z}f(n)=\sum_{m\in\mathbb Z}\widehat f(m).
\]
For $A>0$ the transforms of $e^{-\pi Ax^2}$ and
$xe^{-\pi Ax^2}$ are, respectively,
\[
 A^{-1/2}e^{-\pi y^2/A},\qquad
 -iA^{-3/2}y e^{-\pi y^2/A}.
\]
\end{lemma}
\begin{proof}
The periodization $F(x)=\sum_{n\in\mathbb Z}f(x+n)$ and all its
derivatives converge uniformly for $x\in[0,1]$. Its $m$th Fourier
coefficient is, by an absolutely convergent interchange,
\[
 \int_0^1F(x)e^{-2\pi imx}\,dx
  =\int_{\mathbb R}f(x)e^{-2\pi imx}\,dx=\widehat f(m).
\]
Integration by parts twice makes these coefficients absolutely summable.
The Fourier series of the smooth periodic function $F$ therefore
converges uniformly to $F$. Evaluating at zero proves the sum identity.

Let $G(y)$ be the transform of $e^{-\pi Ax^2}$. Differentiation under
the integral and integration by parts give
$G'(y)=-(2\pi y/A)G(y)$. The Gaussian integral, obtained by squaring
the integral and using polar coordinates, gives $G(0)=A^{-1/2}$.
Solving this differential equation proves the first formula. The transform
of $xe^{-\pi Ax^2}$ is $iG'(y)/(2\pi)$, proving the second.
All differentiations and boundary terms are justified by Gaussian decay.
\end{proof}

Applying the sum identity to $x\mapsto f(r+qx)$ gives, for $q>0$,
\begin{equation}\label{tools:poisson-progression}
 \sum_{n\in\mathbb Z}f(r+qn)
 =\frac1q\sum_{m\in\mathbb Z}
          e^{2\pi imr/q}\widehat f(m/q).
\end{equation}
Indeed, substitution in the Fourier integral gives its $m$th
coefficient as $q^{-1}e^{2\pi imr/q}\widehat f(m/q)$.

\subsection{The gamma function and its zeros}

For $\Re z>0$ define
\[
 \Gamma(z)=\int_0^\infty e^{-t}t^{z-1}\,dt.
\]
The integral and its derivatives converge locally uniformly in this
half-plane. Euler's constant is
$\gamma=\lim_{n\to\infty}(\sum_{k=1}^n k^{-1}-\log n)$.
The limit exists because the sequence is decreasing and positive, as
comparison with the integral of $1/x$ shows.

\begin{lemma}\label{tools:gamma}
The gamma function is nonzero for $\Re z>0$. Its reciprocal extends
to the entire function
\begin{equation}\label{tools:gamma-product}
 \frac1{\Gamma(z)}
   =ze^{\gamma z}\prod_{k=1}^{\infty}
        \left(1+\frac zk\right)e^{-z/k}.
\end{equation}
Its zeros are precisely $0,-1,-2,\ldots$, each simple.
Also $\Gamma(z+1)=z\Gamma(z)$ where both sides are defined, and
$\Gamma(z)\sim1/z$ as $z\to0$.
\end{lemma}
\begin{proof}
For $\Re z>0$, repeated integration by parts gives
\[
 \int_0^1 t^{z-1}(1-t)^n\,dt
  =\frac{n!}{z(z+1)\cdots(z+n)}.
\]
After multiplying by $n^z$ and putting $u=nt$, dominated convergence
applies because $(1-u/n)^n\leq e^{-u}$ for $0\leq u\leq n$.
It follows that
\[
 F_n(z):=\frac{n!n^z}{z(z+1)\cdots(z+n)}\longrightarrow\Gamma(z).
\]
On the other hand, the exact reciprocal is
\[
 F_n(z)^{-1}
  =z\exp\left(z\left(\sum_{k=1}^n\frac1k-\log n\right)\right)
       \prod_{k=1}^n(1+z/k)e^{-z/k}.
\]
On a compact set, the logarithms of the tail factors are uniformly
$O(k^{-2})$. The infinite product therefore converges locally uniformly;
off the nonpositive integers its tail and its finitely many initial
factors are nonzero. Thus the displayed reciprocals converge to a
nonzero function for $\Re z>0$. Passing to the limit in
$F_n(z)F_n(z)^{-1}=1$ proves both nonvanishing of $\Gamma$ and the
product formula there. The product itself is entire. At a nonpositive
integer exactly one factor vanishes simply, which proves the assertion
about its zeros and defines the meromorphic extension of $\Gamma$.
Finally, integration by parts in Euler's integral gives the recurrence.
Since $\Gamma(1)=1$, it also gives the asymptotic at zero.
\end{proof}
